\documentclass[11pt]{article}

\usepackage[margin=1.05in]{geometry}
\usepackage{amsmath,amssymb,amsthm,mathtools}
\usepackage{enumitem}
\usepackage{xcolor}
\usepackage[hypertexnames=false]{hyperref}
\usepackage{microtype}

\hypersetup{
  colorlinks=true,
  linkcolor=blue!55!black,
  urlcolor=blue!55!black,
  citecolor=blue!55!black
}

\setlength{\parindent}{0pt}
\setlength{\parskip}{6pt}
\setlist[itemize]{leftmargin=1.6em,itemsep=2pt,topsep=3pt}
\setlist[enumerate]{leftmargin=1.6em,itemsep=2pt,topsep=3pt}

\newtheorem{theorem}{Theorem}[section]
\newtheorem{proposition}[theorem]{Proposition}
\newtheorem{lemma}[theorem]{Lemma}
\newtheorem{corollary}[theorem]{Corollary}
\newtheorem{assumption}[theorem]{Input}
\newtheorem{remark}[theorem]{Remark}
\newtheorem{definition}[theorem]{Definition}

\newcommand{\R}{\mathbb R}
\newcommand{\E}{\mathbb E}
\newcommand{\ip}[2]{\left\langle #1,#2\right\rangle}
\newcommand{\norm}[1]{\left\|#1\right\|}
\newcommand{\one}{\mathbf 1}
\newcommand{\cH}{\mathcal H}
\newcommand{\He}{\operatorname{He}}
\newcommand{\sgn}{\operatorname{sign}}
\newcommand{\supp}{\operatorname{supp}}
\newcommand{\Proj}{\operatorname{Proj}}
\newcommand{\Var}{\operatorname{Var}}
\newcommand{\Id}{\operatorname{Id}}
\newcommand{\PhiG}{\Phi}

\title{The lower bound proof in 2D}
\author{}
% \date{}

\begin{document}
\maketitle

\begin{abstract}
\noindent This is an expository note on how the proof of the lower bound
\[
        K_G\ge \frac{6\pi}{11}
\]
works. This bound has to be proved for \textbf{all} schemes in any dimension. However, it is first instructive to see how the proof works by specializing in the simplest nontrivial setting: two dimensional schemes.  The point of this note is not that the true theorem is two-dimensional, but that the general case is the same proof with minor changes. 
\end{abstract}

\tableofcontents

\section*{Proof idea}
\addcontentsline{toc}{section}{Proof idea}

\begin{enumerate}
    \item Naor and Regev show that Krivine rounding schemes are optimal: That is, if we can show that over all high-dimensional functions, the value of Krivine rounding scheme value.  For every $f,g$, we have some $\gamma(f,g)$ that we get and we know the upper bound $K_G \leq \pi/2\gamma(f,g)$. Naor and Regev should say the following: $K_G > (1-O(1)/k) \pi/2\gamma(f,g)$ ($k$ is the dimension of the scheme). NEED TO CHECK THIS. 
    \item Claim 1: The Taylor coefficients $b_1, b_3$ of $H_{f,g}(t)$ satisfy $b_3 \geq 2 b_1 - 11/6$. 
    \item Claim 2: From claim 1 alone, we get that $\gamma(f,g) \leq 11/12$. 
\end{enumerate}

 Recall that for a rounding scheme $(f, g)$, we have \[H_{f,g}(t) = b_1 t + b_3 t^3 + b_5 t^5 + \ldots\]
We are going to prove a linear tradeoff between $b_3$ and $b_1$, so that we cannot significantly change one without changing the other. 
\[
        b_3 \ge 2b_1 -\frac{11}{6}.
\]
The fact that this is linear matters because the inequality is valid even for mixed schemes (where you merely average the coefficients of $H$). We call this the \textbf{Affine Strip}. 

\section{Preliminary facts about hermites and gaussians}

Let
\[
        (S,T)\sim N(0,I_2)
\]
with $S,T$ independent standard Gaussians. We define the \textit{inner product} as the Gaussian expectation:
\[
        \ip{F}{G}=\E[F(S,T)G(S,T)].
\]
We use the orthonormal 1D Hermites
\[
        \psi_0=1,\qquad
        \psi_1(s)=s,\qquad
        \psi_2(s)=\frac{s^2-1}{\sqrt2},\qquad
        \psi_3(s)=\frac{s^3-3s}{\sqrt6}.
\]
Any square integrable function has a hermite expansion, written as 
\[F = P_1  F + P_2 F + P_3 F + \ldots \]
where $P_i F$ are projections of $F$ on the $i$-th hermite polynomial $\psi_i$. In our case, we only care about odd functions, so the even hermites of the expansion disappear. 

In two dimensions, the hermites span
\[
        \cH_1=\operatorname{span}\{\psi_1(S),\psi_1(T)\} = \operatorname{span}(S, T).
\]
Here, 
\[
        P_1F=aS+bT,
        \qquad
        a=\E[F(S,T)S],\quad b=\E[F(S,T)T].
\]
Its norm is
\[
        \norm{P_1F}^2=a^2+b^2.
\]

Similarly, the total degree-three hermite span is
\[
        \cH_3=\operatorname{span}\{\psi_3(S),\ \psi_2(S)\psi_1(T),\ \psi_1(S)\psi_2(T),\ \psi_3(T)\}.
\]
So $P_3F$ is the projection onto this four-dimensional space.  If
\[
        P_3F
        =c_{30}\psi_3(S)
         +c_{21}\psi_2(S)\psi_1(T)
         +c_{12}\psi_1(S)\psi_2(T)
         +c_{03}\psi_3(T),
\]
then
\[
        \norm{P_3F}^2=c_{30}^2+c_{21}^2+c_{12}^2+c_{03}^2.
\]


For signs $f,g:\R^2\to\{\pm1\}$, if
\[
        H_{f,g}(t)=\frac\pi2\E[f(X)g(Y_t)]
\]
with $(X,Y_t)$ a pair of correlated standard Gaussian vectors, then Mehler's identity gives
\[
        H_{f,g}(t)=\frac\pi2\sum_{m\ge0}\ip{P_mf}{P_mg}t^m,
\]
Here $b_1=[t]H_{f,g} = \frac\pi2\ip{P_1f}{P_1g}$ and $b_3=[t^3]H_{f,g} = \frac\pi2\ip{P_3f}{P_3g}$.

For the rest of the paper, we will use the shorthand
\[
\nu := \sqrt{\frac{2}{\pi}}.
\]

We want to prove 
\begin{equation}
        b_3 \ge 2b_1 -\frac{11}{6}.
        \tag{Strip}
\end{equation}

For example, the hyperplane partitions lie on this line.  For $f=g=\sgn(S)$,
\[
        \norm{P_1f}=\nu,
        \qquad
        \norm{P_3f}^2=\frac{\nu^2}{6}.
\]
Thus
\[
        (b_1,b_3)=\left(1,\frac16\right),
\]
and this point lies exactly on the line since $1/6=2-11/6$.

Instead of proving the strip directly, it suffices to prove
\begin{equation}
        B_1(f,g)
        :=\norm{P_1f}\norm{P_1g}
          +\ip{P_1f}{P_1g}
          -\ip{P_3f}{P_3g}
        \le \frac{11}{6}\nu^2.
        \tag{A1}
\end{equation}
Indeed, by Cauchy--Schwarz,
\[
        \norm{P_1f}\norm{P_1g}\ge \ip{P_1f}{P_1g}.
\]
So (A1) implies
\[
        2\ip{P_1f}{P_1g}-\ip{P_3f}{P_3g}\le \frac{11}{6}\nu^2.
\]
Multiplying by $\pi/2$ and using $\nu^2=2/\pi$ gives (Strip).

\section{Change of Variables: Agreement and disagreement}

We define a change of variables based on where $f$ and $g$ agree or disagree. Set the functions $h, k$ to be
\[
        h=\frac{f+g}{2},\qquad k=\frac{f-g}{2}.
\]
Then
\[
        f=h+k,
        \qquad
        g=h-k.
\]
Since $f,g$ are signs,
\[
        h,k\in\{-1,0,1\},
        \qquad
        hk=0,
        \qquad
        h^2+k^2=1.
\]
Think of $h$ as the agreement part and $k$ as the disagreement part between $f$ and $g$. This change of variables considerably simplifies the proof, and allows us to reduce the problem to a finite dimensional problem. 

Let
\[
        u=P_1h,\qquad v=P_1k.
\]
Then
\[
        P_1f=u+v,
        \qquad
        P_1g=u-v.
\]
Hence
\[
        \ip{P_1f}{P_1g}=\norm{u}^2-\norm{v}^2.
\]
Also
\[
        \norm{P_1f}\norm{P_1g}
        =\norm{u+v}\norm{u-v}
        \le \norm{u}^2+\norm{v}^2,
\]
using $ab\le(a^2+b^2)/2$ after expanding $\norm{u\pm v}^2$.

For degree three,
\[
        P_3f=P_3h+P_3k,
        \qquad
        P_3g=P_3h-P_3k,
\]
so
\[
        \ip{P_3f}{P_3g}=\norm{P_3h}^2-\norm{P_3k}^2.
\]
Putting these together,
\begin{align*}
        B_1(f,g)
        &\le
        \bigl(\norm{u}^2+\norm{v}^2\bigr)
        +\bigl(\norm{u}^2-\norm{v}^2\bigr)
        -\bigl(\norm{P_3h}^2-\norm{P_3k}^2\bigr)\\
        &=2\norm{P_1h}^2-\norm{P_3h}^2+\norm{P_3k}^2.
\end{align*}

So it is enough to prove
\begin{equation}
\boxed{\displaystyle
        D_1(h,k):=2\norm{P_1h}^2-\norm{P_3h}^2+\norm{P_3k}^2
        \le \frac{11}{6}\nu^2
}
\tag{HK1}
\end{equation}

In order to prove this, we have to handle each of the three expressions $\norm{P_1h}^2, \norm{P_3h}^2, \norm{P_3k}^2$ carefully. In particular, we need to prove upper bounds on  $\norm{P_1h}^2, \norm{P_3k}^2$ and lower bound on $\norm{P_3h}^2$. 

\begin{remark}
    Note that in the inequality before (HK1) the term $\norm{v}^2 = \norm{P_1k}^2$ conveniently cancels. This is a result of our choice of $11\pi/6$, which leads to the equation $b_3 \geq 2b_1 - \frac{11}{6}$, where the term $2b_1$ gives rise to the cancellation. We can prove stronger lower bounds with the same method in the paper, but this is more complicated because we in addition have to bound $\norm{v}^2 = \norm{P_1k}^2$.
\end{remark}

\section{Change of Variable: Rotation}

The linear part of $h$ has the form
\[
        P_1h(S,T)=aS+bT.
\]
The coefficient vector is $(a,b)$, and
\[
        \norm{P_1h}=\sqrt{a^2+b^2}.
\]
If $(a,b)\neq0$, rotate coordinates so the first coordinate points in the $(a,b)$ direction. After this rotation, and after renaming the coordinates, we have
\[
        P_1h(S,T)=\norm{P_1h}\,S.
\]
If $P_1h=0$, choose any direction as $S$.

Define
\[
        x:=\frac{\norm{P_1h}}{\nu}\in[0,1].
\]
Let $Z \sim N(0,1)$, so that $\mathbb{E}|Z| = \nu$. The bound $x\le1$ follows from
\[
        \norm{P_1h}=\sup_{\alpha^2+\beta^2=1}\E[h(S,T)(\alpha S+\beta T)]
        \le \E|Z|=\nu.
\]
After rotation,
\[
        \ip{h}{S}=\norm{P_1h}=\nu x.
\]

\section{The agreement part}

The agreement part of $D_1$ is $2\norm{P_1h}^2-\norm{P_3h}^2$.

The helpful negative term is
\[
        -\norm{P_3h}^2.
\]
To use it, we only need a lower bound on $\norm{P_3h}^2$.

In two dimensions,
\[
        \norm{P_3h}^2
        \ge \ip{h}{\psi_3(S)}^2,
\]
because $\psi_3(S)$ is one of the four orthonormal degree-three basis elements.

The trick we will use is to condition on one variable, so define
\[
        A(s):=\E_T[h(s,T)].
\]
This is the average value of $h$ along the horizontal slice $S=s$.  Since $h\in[-1,1]$,
\[
        |A(s)|\le1.
\]
Also, for every function $\varphi$ depending only on $S$, it is a straightforward conditioning to show that 
\[
        \E[h(S,T)\varphi(S)]=\E[A(S)\varphi(S)].
\]
In particular,
\[
        \E[A(S)S]=\ip{h}{S}=\nu x
\]
and
\[
        \ip{h}{\psi_3(S)}=\E[A(S)\psi_3(S)].
\]

This is the clever trick. Note that the problem is now a variational problem in the one variable function $A$, even though we were originally working in two dimensions. Now in order to bound the agreement part, it is sufficient to relax $A$, and answer the following question

\begin{quote}
``Among all $A:\R\to[-1,1]$ with $\E[A(S)S]=\nu x$, how small can the cubic coefficient be?''
\end{quote}

The answer is the rearrangement lemma.

Define
\[
        W(S):=S^3-3S=\sqrt6\,\psi_3(S).
\]
Also define
\[
        \Delta(x):=x+2(1+x)\log\frac{1+x}{2},
        \qquad
        \Delta_+(x):=\max\{\Delta(x),0\}.
\]

\begin{lemma}[One-dimensional rearrangement]
Let $S\sim N(0,1)$ and $x\in[0,1]$.  Then
\[
        \sup_{\substack{|A|\le1\\ \E[A(S)S]=\nu x}}
        \E[A(S)(S^3-3S)]
        =-\nu\Delta(x).
\]
\end{lemma}

We delegate the proof to Appendix~\ref{app:rearrangement_lemma}, which uses standard variational techniques. Using this rearrangement lemma then gives us the desired bound on the agreement part, that we state below. 

\begin{corollary}[Diagonal bound]
For the two-dimensional ternary $h$ above,
\[
        \norm{P_3h}^2\ge \frac{\nu^2}{6}\Delta_+(x)^2.
\]
Consequently
\[
        2\norm{P_1h}^2-\norm{P_3h}^2
        \le
        \nu^2\left(2x^2-\frac16\Delta_+(x)^2\right).
\]
\end{corollary}

\section{The disagreement part}

Now we have to bound the disagreement term
\[
        \norm{P_3k}^2.
\]

Define the one-dimensional fibers (think "slices" across one variable) 
\[
        u_t(s):=k(s,t).
\]
Since $k\in\{-1,0,1\}$,
\[
        u_t:\R\to\{-1,0,1\}.
\]

For $j=0,1,2,3$, define
\[
        K_j(t):=\E_S[k(S,t)\psi_j(S)].
\]
Thus $K_j$ is the $j$th Hermite coefficient of the $S$-fiber, as a function of $t$.

The four degree-three basis elements are
\[
        \psi_3(S),\qquad
        \psi_2(S)\psi_1(T),\qquad
        \psi_1(S)\psi_2(T),\qquad
        \psi_3(T).
\]
Their coefficients in $k$ are, respectively,
\[
        \E_T[K_3(T)],
\]
\[
        \E_T[K_2(T)\psi_1(T)],
\]
\[
        \E_T[K_1(T)\psi_2(T)],
\]
\[
        \E_T[K_0(T)\psi_3(T)].
\]
Equivalently,
\begin{equation}
        \norm{P_3k}^2
        =\norm{\Pi_0K_3}^2+
          \norm{\Pi_1K_2}^2+
          \norm{\Pi_2K_1}^2+
          \norm{\Pi_3K_0}^2,
        \tag{3}
\end{equation}
where $\Pi_j$ is the one-dimensional projection in the $T$ variable onto Hermite degree $j$. Since $\Pi_j$ are orthogonal projections onto a closed subspace, it decreases norm, so we get the simpler upper bound 
\[
        \norm{P_3k}^2
        \le \sum_{j=0}^3\norm{K_j}^2
        =\E_T\sum_{j=0}^3\left(\E_S[k(S,T)\psi_j(S)]\right)^2.
\]
For a one dimensional ternary function $u$, define
\[
        V(u):=\sum_{j=0}^3\left(\E[u(Z)\psi_j(Z)]\right)^2.
\]
Then the previous inequality becomes
\begin{equation}
        \boxed{\norm{P_3k}^2\le \E_T V(u_T).}
        \tag{4}
\end{equation}\[p_a(Z) = \sum\limits_{j=0}^3 a_j \psi_j(Z), \qquad ||a||_2 = 1\] 

For a ternary function $u:\R\to\{-1,0,1\}$, define its weighted support budget
\[
        \beta(u):=\E[|Z|u(Z)^2].
\]
Since $u^2=\one_{u\neq0}$,
\[
        \beta(u)=\E[|Z|\one_{u(Z)\neq0}].
\]
So $\beta(u)$ is the $|Z|$-weighted Gaussian size of the support of the fiber.

\begin{theorem}[Fiber inequality]
For every ternary $u:\R\to\{-1,0,1\}$,
\begin{equation}
        V(u)\le 3\nu\beta(u)-\beta(u)^2.
        \tag{FQ1}
\end{equation}
\end{theorem}

\begin{proof}
    We are optimizing over infinitely many $u$ functions, and our approach of dealing with this by eliminating $u$ entirely. We will then verify the remaining steps in Arb. 

    Recall that the Euclidean norm has a dual, that is, one can write  $||x||_2 = \sup\limits_{\lVert a \rVert_2 = 1} a \cdot x$. Then, write $V(u) = \sum_{j=0}^3 m_i^2$, where $m_i = \E[u(Z)\psi_j(Z)]$. Using the dual form, we get \[\sqrt{V(u)} = \sup\limits_{\lVert a \rVert_2 = 1} \sum\limits_{j=0}^3 a_j m_j.\]
    The nice thing about this form is that it linearizes the expression. Since \[\sum\limits_{j=0}^3 a_j m_j =\mathbb{E}\left[u(Z) \sum\limits_{j=0}^3 a_j \psi_j(Z) \right],\]
    define the cubic polynomial
    \[p_a(Z) = \sum\limits_{j=0}^3 a_j \psi_j(Z), \qquad \lVert a \rVert_2 = 1,\] which gives \[\sqrt{V(u)} = \sup\limits_{\lVert a \rVert_2 = 1} \mathbb{E}[u(Z) p_a(Z)].\]
    Recall here one more time that $u$ is a ternary function taking values in $\{-1, 0, 1\}$, so at each point $Z=z_0$, the best value of $u$ to take is $\sgn(p_a(z_0))$. However, recall that we cannot do this arbitrarily, since we have a weighted budget constraint of $\beta(u) = \mathbb{E}[|Z| u(Z)^2]$. So, any time we set $u$ to be $\sgn(p_a(z_0))$, we gain $|p_a(Z)|$ in the expectation but lose $|z|$ in the budget. Our goal is to find a set $E$ of points maximizing our objective, so that for any point $z_0 \in E$ ,we choose $u$ to be $\sgn(p_a(z_0))$, and otherwise set $u$ to $0$.    

    By a standard Lagrange multipliers argument, the right set here is to  choose a threshold value $\tau$ and take points in the set $E = \{z : |p(z)| \ge \tau  |z|\}$. Then for each $a$, the best function $u$ is \[u(z) = \begin{cases}
        \sgn(p(z_0)) & \text{if } |p(z)| \ge \tau |z| \\
        0 & \text{else} \\  
    \end{cases}\]
    Thus for each $a$, the best $u$ is fixed! We can only worry about $a$  and $\tau$ now. Let $m = \mathbb{E}[|p_a(Z)| \mathbb{1}_{E}]$ and $\beta$ the resulting budget $\mathbb{E}[|Z| 1_{E}]$. The required inequality to prove is \[m^2 \le 3 \nu \beta - \beta^2.\]
    This is now a problem over finite dimensional parameters $a$ and $\tau$. That is a numerical problem that can be checked with Arb. Below we explain how. 

For fixed \(a,\tau\), the set
\[
E_{a,\tau}
=
{|p_a(z)|\ge \tau |z|}
\]

is defined by a one-dimensional polynomial inequality. Since \(p_a\) is cubic, the boundary points come from equations like
\[
p_a(z)=\tau z
\]
or
\[
p_a(z)=-\tau z,
\]
with sign changes around \(0\). These are polynomial root problems in one variable.

Once the roots are enclosed, \(E_{a,\tau}\) is a union of intervals. Then the needed integrals are Gaussian moments over intervals:

\[
\int_I |p_a(z)|,d\gamma(z),
\qquad
\int_I |z|,d\gamma(z),
\qquad
\int_I \psi_j(z),d\gamma(z).
\]

These can be computed rigorously with Arb.    
\end{proof}

 

Apply (FQ1) to each fiber $u_T$.  From (4),
\[
        \norm{P_3k}^2
        \le
        \E_T\left[3\nu\beta(u_T)-\beta(u_T)^2\right].
\]
We now prove the following bound. Note that this step reduces the dimension down to 1.

\begin{theorem}[Fiber Residual Bound \#1]
We have the fiber bound 
\begin{equation}
        \norm{P_3k}^2
        \le
        \nu^2\left(3(1-x)-(1-x)^2\right).
\end{equation}
\end{theorem}

\begin{proof}
We do this by relating the average budget to $x$.
Since
\[
        \nu x=\E[hS]
\]
and
\[
        hS\le |S|h^2,
\]
we get
\begin{align*}
        \nu x
        &\le \E[|S|h^2]\\
        &=\E[|S|(1-k^2)]\\
        &=\nu-\E[|S|k^2].
\end{align*}
Therefore
\[
        \E[|S|k^2]\le \nu(1-x).
\]
But
\[
        \E[|S|k^2]
        =\E_T\E_S[|S|k(S,T)^2]
        =\E_T\beta(u_T).
\]
Let
\[
        b:=\E_T\beta(u_T).
\]
Then
\[
        b\le \nu(1-x).
\]

The function
\[
        \phi(t)=3\nu t-t^2
\]
is concave, so Jensen gives
\[
        \E_T\phi(\beta(u_T))\le \phi(\E_T\beta(u_T))=3\nu b-b^2.
\]
On $[0,\nu]$, $\phi$ is increasing.  Hence, using $b\le \nu(1-x)$,
\begin{equation}
        \norm{P_3k}^2
        \le
        \nu^2\left(3(1-x)-(1-x)^2\right).
        \tag{5}
\end{equation}
\end{proof}


This is our first fiber residual bound. But this bound is not best for small $x$.  We also use a second cruder bound. 

\begin{theorem}[Fiber Residual Bound \#2]
    \begin{equation}
        \norm{P_3k}^2
        \le
        2\PhiG(\sqrt{-2\log x})-1.
        \tag{6}
\end{equation}
\end{theorem}

\begin{proof}

Let
\[
        \rho(s):=\E_T[k(s,T)^2].
\]
Then $0\le\rho\le1$ and
\[
        \norm{k}^2_2=\E\rho(S).
\]
Also
\[
        \E[|S|\rho(S)]=\E[|S|k(S,T)^2]\le \nu(1-x).
\]
Since $P_3$ is a projection,
\[
        \norm{P_3k}^2\le \norm{k}^2_2=\E\rho(S).
\]

Given a budget on $\E[|S|\rho(S)]$, the way to maximize $\E\rho(S)$ is to spend the budget where $|S|$ is smallest.  Thus the extremal support is centered around zero.

Let
\[
        r_x:=\sqrt{-2\log x}
\]
for $x\in(0,1]$, with the convention $r_0=\infty$.  Then
\[
        \E[|S|\one_{|S|\le r_x}]=\nu(1-e^{-r_x^2/2})=\nu(1-x).
\]
Therefore the bathtub principle gives
\begin{equation}
        \norm{P_3k}^2
        \le
        \mathbb P(|S|\le r_x)
        =2\PhiG(\sqrt{-2\log x})-1.
        \tag{6}
\end{equation}
\end{proof}

Combining the two fiber residual bounds, define
\begin{equation}
        R(x):=
        \min\left\{
        3(1-x)-(1-x)^2,
        \frac{1}{\nu^2}\left(2\PhiG(\sqrt{-2\log x})-1\right)
        \right\}.
        \tag{7}
\end{equation}
Then
\[
        \norm{P_3k}^2\le \nu^2R(x).
\]

\section{Putting it together}

From the diagonal bound and the residual bound,
\[
        D_1(h,k)
        \le
        \nu^2\left(2x^2-\frac16\Delta_+(x)^2+R(x)\right).
\]
So it remains to prove the scalar inequality
\begin{equation}
        E_1(x):=2x^2-\frac16\Delta_+(x)^2+R(x)
        \le \frac{11}{6}
        \qquad(0\le x\le1).
        \tag{8}
\end{equation}

Let
\[
        r_-:=\frac{1-1/\sqrt3}{2},
        \qquad
        r_+:=\frac{1+1/\sqrt3}{2}.
\]
These are the roots of
\[
        x^2-x+\frac16=0.
\]

Now we do some simple casework based on the range in which $x$ lies. Everything is simple algebra.  

\subsection*{Low range: $0\le x\le r_-$}

Use the crude bound $R(x)\le 1/\nu^2=\pi/2$ and drop the negative term:
\[
        E_1(x)\le 2x^2+\frac\pi2
        \le 2r_-^2+\frac\pi2.
\]
Now
\[
        2r_-^2+\frac\pi2
        =\frac23-\frac1{\sqrt3}+\frac\pi2
        <\frac{11}{6}.
\]

\subsection*{Middle range: $r_-\le x\le r_+$}

Use the fiber branch of $R$ and again drop the negative term:
\begin{align*}
        E_1(x)
        &\le 2x^2+3(1-x)-(1-x)^2\\
        &=x^2-x+2.
\end{align*}
Since $x\in[r_-,r_+]$,
\[
        x^2-x+\frac16\le0.
\]
Therefore
\[
        x^2-x+2\le \frac{11}{6}.
\]

\subsection*{High range: $r_+\le x\le1$}

Use the fiber branch again.  We need the cubic correction.
It suffices to show
\begin{equation}
        \Delta(x)^2\ge 6x^2-6x+1.
        \tag{9}
\end{equation}
Indeed, then
\begin{align*}
        E_1(x)
        &\le 2x^2+3(1-x)-(1-x)^2-\frac16(6x^2-6x+1)\\
        &=\frac{11}{6}.
\end{align*}

To prove (9), first note that
\[
        \Delta(x)-(3x-2)
        =2\left((1-x)+(1+x)\log\frac{1+x}{2}\right).
\]
Writing $t=(1+x)/2$, the bracket becomes
\[
        2\bigl((1-t)+t\log t\bigr)\ge0
\]
for $0<t\le1$.  Hence
\[
        \Delta(x)\ge 3x-2.
\]
For $x\ge r_+$, both sides needed below are nonnegative, and
\[
        (3x-2)^2-(6x^2-6x+1)=3(1-x)^2\ge0.
\]
Thus
\[
        \Delta(x)^2\ge (3x-2)^2\ge 6x^2-6x+1.
\]
This proves the high range.

Combining the three ranges proves (8).

We have thus shown
\[
        D_1(h,k)\le \frac{11}{6}\nu^2.
\]
As explained earlier, this implies the affine inequality
\[
        b_3 \ge 2b_1 -\frac{11}{6}.
\]
Since this is affine, it is preserved under mixtures. 

\section{How the strip gives $6\pi/11$}

Here we briefly mention the analytic nuisances that need to be cleared up in order to convert everything into a proof. 

\begin{assumption}[Standard closure/transfer lemma]
If an affine coefficient inequality
\[
        b_3\ge 2b_1-\frac{11}{6}
\]
holds for all atom kernels in every dimension, then it also holds for the ideal spherical kernel $c/K$ for every $K>K_G$, where
\[
        c(\theta,\phi)=\ip{\theta}{\phi}.
\]
\end{assumption}

The reason this is true is that, for $K>K_G$, the kernel $c/K$ lies in the weak-* closed convex hull of sign atom kernels.  Since the strip is affine and weak-* closed, it survives the passage to this limit.

For the ideal kernel $c/K$, the first coefficient tends to
\[
        b_1(c/K)\longrightarrow \frac{\pi}{2K}
\]
as the sphere dimension goes to infinity, while the third coefficient tends to
\[
        b_3(c/K)\longrightarrow 0.
\]
Applying the strip gives
\[
        0\ge 2\cdot\frac{\pi}{2K}-\frac{11}{6}
        =\frac\pi K-\frac{11}{6}.
\]
Thus
\[
        K\ge \frac{6\pi}{11}.
\]
Since this holds for every $K>K_G$,
\[
        K_G\ge \frac{6\pi}{11}.
\]


\section{What changes in higher dimensions}

In higher dimensions, write the Gaussian input as
\[
        X=(S,Y),
\]
where $S$ is the spine coordinate and $Y$ is the transverse Gaussian vector.

Nothing changes in the agreement part.  Define
\[
        A(s):=\E_Y[h(s,Y)].
\]
Then
\[
        \E[A(S)S]=\nu x,
        \qquad
        \ip{h}{\psi_3(S)}=\E[A(S)\psi_3(S)].
\]
The same rearrangement lemma gives the same lower bound on $\norm{P_3h}^2$.

For the disagreement part, the only change is notation.  In two dimensions, total degree three decomposed as
\[
        3=3+0=2+1=1+2=0+3.
\]
In general this becomes the orthogonal decomposition
\[
        \cH_3(S,Y)
        =\bigoplus_{j=0}^3 \cH_j(S)\otimes \cH_{3-j}(Y).
\]
For each fiber define
\[
        u_Y(s):=k(s,Y),
        \qquad
        K_j(Y):=\E_S[k(S,Y)\psi_j(S)].
\]
Then
\[
        P_3k
        =\sum_{j=0}^3 \psi_j(S)\,\Pi_{3-j}^Y K_j(Y),
\]
where $\Pi_{3-j}^Y$ projects onto transverse Hermite degree $3-j$.  Therefore
\[
        \norm{P_3k}^2
        =\sum_{j=0}^3\norm{\Pi_{3-j}^Y K_j}^2
        \le \sum_{j=0}^3\norm{K_j}^2
        =\E_Y V(u_Y).
\]
This is the exact higher-dimensional version of the two-dimensional calculation.

The budget calculation is also identical:
\[
        \E_Y\beta(u_Y)=\E[|S|k(S,Y)^2]\le \nu(1-x).
\]
Then the same fiber inequality, Jensen step, bathtub bound, and scalar inequality finish the proof.

\appendix 

\section{Proof of rearrangement lemma}
\label{app:rearrangement_lemma}


\begin{lemma}[One-dimensional rearrangement]
Let $S\sim N(0,1)$ and $x\in[0,1]$.  Then
\[
        \sup_{\substack{|A|\le1\\ \E[A(S)S]=\nu x}}
        \E[A(S)(S^3-3S)]
        =-\nu\Delta(x).
\]
\end{lemma}

\begin{proof}
Fix $r\ge0$ by
\[
        e^{-r^2/2}=\frac{1+x}{2},
\]
and put
\[
        \mu=r^2-3.
\]
Then
\[
        W-\mu S=S^3-3S-(r^2-3)S=S(S^2-r^2).
\]

For any feasible $A$,
\begin{align*}
        \E[AW]
        &=\E[A(W-\mu S)]+\mu\E[AS]\\
        &=\E[A(W-\mu S)]+\mu\nu x\\
        &\le \E|W-\mu S|+\mu\nu x,
\end{align*}
using $|A|\le1$.

Equality is attained by
\[
        A_*(S)=\sgn(W-\mu S)=\sgn(S)\sgn(S^2-r^2),
\]
provided $A_*$ has the right first moment.  It does, because
\[
        A_*(S)S=|S|\sgn(S^2-r^2),
\]
so
\begin{align*}
        \E[A_*(S)S]
        &=2\E[|S|\one_{|S|>r}]-\E|S|\\
        &=2\nu e^{-r^2/2}-\nu\\
        &=\nu x.
\end{align*}
Thus $A_*$ is feasible and attains the dual upper bound.

It remains to compute its value.  Since
\[
        \sgn(S)(S^3-3S)=|S|^3-3|S|,
\]
we get
\begin{align*}
        \E[A_*W]
        &=\E[(|S|^3-3|S|)\sgn(S^2-r^2)]\\
        &=2\E[(|S|^3-3|S|)\one_{|S|>r}]
          -\E[|S|^3-3|S|].
\end{align*}
The Gaussian tail moments are
\[
        \E[|S|\one_{|S|>r}]=\nu e^{-r^2/2},
        \qquad
        \E[|S|^3\one_{|S|>r}]=\nu(r^2+2)e^{-r^2/2}.
\]
Also
\[
        \E|S|=\nu,
        \qquad
        \E|S|^3=2\nu.
\]
Therefore
\begin{align*}
        \E[A_*W]
        &=2\nu(r^2-1)e^{-r^2/2}+\nu\\
        &=\nu\left(1+(1+x)(r^2-1)\right)\\
        &=\nu\left((1+x)r^2-x\right).
\end{align*}
Since $r^2=-2\log((1+x)/2)$,
\[
        \E[A_*W]
        =-\nu\left(x+2(1+x)\log\frac{1+x}{2}\right)
        =-\nu\Delta(x).
\]
This proves the lemma.
\end{proof}

\end{document}